\section*{Capítulo \ref{ch:b3:glaciology} --- Glaciología}

\begin{solution}{exo:b3:glaciology:1}
$\tau_b = \rho gH\sin\alpha = 917\times9.81\times250\times\sin\ang{3} =
\qty{118}{kPa}$, unos \qty{1.2}{bar}: el orden de los esfuerzos bajo
glaciares reales.
\end{solution}

\begin{solution}{exo:b3:glaciology:2}
(a)~$\qty{2.9e15}{m^3}\times917 = \qty{2.66e18}{kg} = \qty{2.66e6}{Gt}$;
dividido por \qty{360}{Gt/mm}, resulta \qty{7390}{mm}, \qty{7.4}{m}, el valor citado.
(b)~$\qty{3.0e16}{m^3}\times917/360 = \qty{76}{m}$, frente a \qty{58}{m}.
El hielo de Groenlandia descansa en su mayor parte sobre lecho rocoso por encima del nivel del mar, así que todo él
añadiría agua al océano. Buena parte del hielo antártico descansa sobre un lecho
situado por debajo del nivel del mar: ese hielo ya desplaza agua marina, y fundirlo solo sustituye
el agua que desplaza; las plataformas de hielo flotantes no añaden nada en absoluto.
\end{solution}

\begin{solution}{exo:b3:glaciology:3}
Ganancia $\num{6e6}\times1.2 = \qty{7.2e6}{m^3}$, pérdida $\num{4e6}\times2.1 =
\qty{8.4e6}{m^3}$ de agua: neto \qty{-1.2e6}{m^3}, un balance específico medio
de $-1.2\times10^6/10^7 = \qty{-0.12}{m}$ equivalente en agua. La
razón de área de acumulación es $6/10 = 0.6$. El glaciar está perdiendo masa:
reduciéndose.
\end{solution}

\begin{solution}{exo:b3:glaciology:4}
$1 - 820/917 = \num{0.106}$ y $1 - 830/917 = \num{0.095}$: alrededor de una décima
parte del volumen sigue siendo espacio poroso cuando se produce el \omterm{def:b3:glaciology:firn}{cierre de los poros}.
\end{solution}

\begin{solution}{exo:b3:glaciology:5}
$\tau_b = 917\times9.81\times300\times\sin\ang{3} = \qty{141.2}{kPa}$. Por
el teorema de la losa, $u_s = \tfrac{2A}{4}\tau_b^3H = 0.5\times\num{1.43e-16}
\times(\num{1.412e5})^3\times300 = \qty{60}{m/yr}$. Con el doble de espesor, $u_s
\propto H^{n+1} = H^4$, dieciséis veces más rápido. Con la mitad de pendiente: $u_s \propto
(\sin\alpha)^3$, ocho veces más lento.
\end{solution}

\begin{solution}{exo:b3:glaciology:6}
A \qty{-25}{\celsius} (\qty{248.15}{K}), $A = \num{3.61e-13}\exp(-60000/
(8.314\times248.15)) = \qty{8.5e-26}{Pa^{-3}.s^{-1}}$, o
\qty{2.7e-18}{Pa^{-3}.yr^{-1}}. A \qty{-10}{\celsius}, $\num{3.61e-13}
\exp(-27.42) = \qty{4.4e-25}{Pa^{-3}.s^{-1}} = \qty{1.4e-17}{Pa^{-3}.yr^{-1}}$.
A \qty{0}{\celsius} (rama cálida) \qty{1.4e-16}{Pa^{-3}.yr^{-1}}. Bajo
los mismos esfuerzos, las velocidades están en la razón de los \omterm{def:b3:glaciology:glen}{factores de velocidad}:
$\num{1.43e-16}/\num{2.67e-18} = 54$.
\end{solution}

\begin{solution}{exo:b3:glaciology:7}
$h = \sqrt{2\times10^5\times10^6/(917\times9.81)} = \qty{4.7}{km}$, próximo a
los \qty{4.9}{km} de la Antártida: el modelo plástico obtiene el espesor de un
casquete de hielo solo a partir de su anchura.
\end{solution}

\begin{solution}{exo:b3:glaciology:8}
$\tau = 200/5 = \num{40}\text{ años}$. La longitud seguiría ajustándose durante
varias décadas, unos pocos tiempos de respuesta, tras cesar el calentamiento: los glaciares actuales
no han alcanzado todavía el clima actual.
\end{solution}

\begin{solution}{exo:b3:glaciology:9}
(a)~$h = \sqrt{2\times2.03\times20\times\num{8.64e6}/(917\times\num{3.34e5})}
= \qty{1.51}{m}$. Como $h^2 \propto t$, \qty{3}{m} tarda $100\times
(3/1.51)^2 = 393$ días, más que un invierno. (b)~El hielo mucho más grueso que
\qty{3}{m} no puede formarse por congelación en un solo invierno: es hielo que ha
sobrevivido a varios inviernos, o hielo amontonado por el viento y las corrientes en crestas.
\end{solution}

\begin{solution}{exo:b3:glaciology:10}
(a)~$H = (1025/917)\times800 = \qty{894}{m}$. (b)~A \qty{900}{m} el
espesor en la \omterm{def:b3:glaciology:grounding-line}{línea de apoyo} es \qty{1006}{m}, $9/8$ del anterior; el flujo
aumenta en $(9/8)^5 = 1.8$. (c)~A medida que la línea retrocede por un lecho que se profundiza
hacia el interior, el hielo en la línea es más grueso, vierte más caudal, adelgaza el hielo
detrás, y la línea retrocede aún más: una retroalimentación positiva. Una plataforma de hielo
que sujeta el hielo apoyado reduce el flujo en la \omterm{def:b3:glaciology:grounding-line}{línea de apoyo} y puede
mantenerla en su sitio incluso sobre una pendiente descendente.
\end{solution}

\begin{solution}{exo:b3:glaciology:11}
Con $h(r) = c\sqrt{R - r}$, $c = \sqrt{2\tau_0/(\rho g)}$, $V =
\int_0^R2\pi rc\sqrt{R - r}\,\mathrm dr$. Sustituyendo $s = R - r$,
\[
\int_0^R(R - s)\sqrt s\,\mathrm ds = \tfrac23R^{5/2} - \tfrac25R^{5/2} =
\tfrac4{15}R^{5/2},
\]
de donde $V = \frac{8\pi}{15}cR^{5/2}$. Al duplicar $R$
se multiplican $V$ y el \omterm{def:b3:glaciology:sle}{equivalente en nivel del mar} por $2^{5/2} = 5.7$.
\end{solution}

\begin{solution}{exo:b3:glaciology:12}
(a)~Marzo: $14.13/16.34 - 1 = \qty{-13.5}{\percent}$; septiembre: $4.81/7.05
- 1 = \qty{-31.8}{\percent}$. (b)~En marzo el borde del hielo se sostiene sobre todo por
las costas que rodean el océano Ártico, y la recongelación cada invierno
vuelve a llenar la cuenca; el hielo más delgado de los últimos años se funde
por completo cada verano. Y la retroalimentación hielo--albedo actúa en verano: el agua libre
absorbe alrededor de \qty{94}{\percent} de la luz solar, el hielo refleja de
\qtyrange{50}{70}{\percent}, por lo que cada pérdida de hielo estival calienta el agua y
funde más hielo.
\end{solution}

\begin{solution}{pb:b3:glaciology:1}
\textbf{1.} La columna sobre el lecho, de peso $\rho gH$ por unidad de área, es
empujada pendiente abajo por la componente $\rho gH\sin\alpha$, equilibrada por el
esfuerzo de cizalla en su base; para pendientes pequeñas $\sin\alpha \approx \alpha$ (el
teorema de la losa con la pendiente superficial).

\textbf{2.} $917\times9.81\times580\times0.0190 = \qty{99.1}{kPa}$;
$917\times9.81\times470\times0.0240 = \qty{101.5}{kPa}$; $917\times9.81
\times335\times0.0330 = \qty{99.4}{kPa}$.

\textbf{3.} El esfuerzo es el mismo dentro de \qty{2}{\percent} mientras el
espesor casi se reduce a la mitad: el casquete se comporta casi como un material plástico, con $\tau_0 \approx
\qty{100}{kPa}$.

\textbf{4.} $2\times10^5/(917\times9.81) = \qty{22.23}{m}$.

\textbf{5.} $h = \sqrt{22.23\,(R - r)}$: \qty{577}{m}, \qty{472}{m} y
\qty{333}{m} a 15, 10 y \qty{5}{km} del margen, a menos de unos pocos metros
de los sondeos.

\textbf{6.} $\sqrt{22.23\times20000} = \qty{667}{m}$.

\textbf{7.} Como en el \cref{exo:b3:glaciology:11}: $V = \int_0^R2\pi r
\sqrt{22.23(R - r)}\,\mathrm dr = \frac{8\pi}{15}\sqrt{2\tau_0/(\rho
g)}\,R^{5/2}$.

\textbf{8.} $\frac{8\pi}{15}\times4.715\times(\num{2e4})^{5/2} =
\qty{4.47e11}{m^3} = \qty{447}{km^3}$.

\textbf{9.} $\pi\times20^2 = \qty{1257}{km^2}$; espesor medio
$447/1257 = \qty{0.356}{km}$, \qty{356}{m}, algo más de la mitad del
espesor central.

\textbf{10.} $2^{5/2} = 5.7$ veces.

\textbf{11.} $A = \num{3.61e-13}\exp(-60000/(8.314\times263.15)) =
\qty{4.45e-25}{Pa^{-3}.s^{-1}}$, por \qty{3.156e7}{s}:
\qty{1.40e-17}{Pa^{-3}.yr^{-1}}.

\textbf{12.} A \qty{10}{km}, $H = \qty{471.5}{m}$, pendiente $471.5/20000 =
\num{0.0236}$, $\tau_b = \qty{100}{kPa}$; $u_s = \tfrac{2A}{4}\tau_b^3H =
0.5\times\num{1.40e-17}\times10^{15}\times471.5 = \qty{3.3}{m/yr}$.

\textbf{13.} $q = \tfrac{2A}{5}\tau_b^3H^2 = 0.4\times\num{1.40e-17}
\times10^{15}\times471.5^2 = \qty{1250}{m^2/yr}$; velocidad media $q/H =
\qty{2.6}{m/yr}$, cuatro quintos de la velocidad superficial.

\textbf{14.} $2\pi rq = \pi r^2\bar b$, luego $\bar b = 2q/r = 2\times1250/10^4
= \qty{0.25}{m}$ de hielo al año sobre los \qty{10}{km} interiores.

\textbf{15.} $h = 450$ donde $R - r = 450^2/22.23 = \qty{9.11}{km}$, en
$r = \qty{10.9}{km}$; el área de acumulación es el disco interior a él, $(10.9/
20)^2 = 0.30$ del casquete.

\textbf{16.} $A$ aumenta \num{10.2} veces, y con ello las velocidades bajo los mismos
esfuerzos. El hielo saldría más deprisa de lo que aporta el balance: el
casquete se adelgazaría y se extendería hasta que los esfuerzos, menores, volvieran a igualar el
balance; un casquete plástico más blando tiene un esfuerzo de fluencia menor y un
perfil más plano.

\textbf{17.} $\num{4.47e11}\times917 = \qty{4.10e14}{kg} = \qty{410}{Gt}$.

\textbf{18.} $410/360 = \qty{1.14}{mm}$.

\textbf{19.} $356/0.5 = 711$ años.

\textbf{20.} A medida que la superficie desciende, queda en aire más cálido y se funde más deprisa,
la retroalimentación altitud--balance de masa; y a medida que se adelgaza, cambian también los esfuerzos y el
flujo. El adelgazamiento se acelera.

\textbf{21.} $410/243 = \num{1.7}\text{ años}$: todo el casquete equivale a menos de dos
años de la pérdida de Groenlandia en 2010--2019.

\textbf{22.} El casquete de hielo de Arkel tiene un \omterm{def:b3:glaciology:sle}{equivalente en nivel del mar} de
\textbf{\qty[text-series-to-math=true]{1.1}{mm}} (\qty{1.14}{mm}).
\end{solution}
